\scanpage{162}
Soit $f:X\mathrel{\unsure{-|\to}}Y$ un morphisme \inserted{propre} dominant, $Y$ \inserted{loc. noeth.} \uline{normal} \uline{connexe}.
Alors $\pi_1(X)\to\pi_1(Y)$ a un conoyau fini (regardant
le pt générique $y$ de $Y$, et notant que $\pi_1(y)\to\pi_1(Y)$ est
surjectif, on est ramené au cas où $Y=\operatorname{Spec}(K)$.
On voit de \unsure{$=$|même} qu'on peut supposer $X$ séparable sur $k$
et \unsure{abst|absolt} connexe [remplacer $K$ par une extension \inserted{séparable}
finie convenable $K'$], et diviser \deleted{\unsure{$O$}} $X'$ par ses élts
nilpotents, enfin prendre une composante connexe
de $X'$]. Mais alors $\pi_1(X)\to\pi_1(Y)$ est \uline{surjectif}..

\noindent\rule{2em}{0.25pt}

N.B. Si $Y$ n'est pas \uline{normal}, le th est faux,
\unsure{$=$|même|si} si p.ex $X$ est le normalisé de $Y$. Ex. $Y$ courbe
\unsure{rat.} sur $k$ alg. clos, avec pt double ordinaire,
$X'$ sa normalisée. On a $\pi_1(X)=0$, $\pi_1(Y)=\widehat{\mathbb Z}$
privé de \deleted{\unsure{sa}} $p$-composante.

\noindent\rule{2em}{0.25pt}

On voit facilement que dans le dernier
\unsure{exemple|exemples}, on ne peut former $\underline{\operatorname{Hom}}_Y(X,Y\amalg Y)$
[car sur un \inserted{\unsure{certain}} pt $y_0$ de $Y$, il y a deux pts
géométriques dans l'image inverse, alors
qu'il n'y en a qu'un au voisinage].

\uline{Supposons cependant} $f:\unsure{X|x}\to Y$ \uline{universellement}
\uline{ouvert} [p. ex. $Y$ \unsure{normale|unibranche}], je crois qu'on
% archive sheet number: 161
